\section*{Purpose and Goals}

\par 
The Free Protocols Foundation (FPF) is a global independent public forum,
dedicated to the support of the internet engineering profession towards 
our profession's responsibility to society and humanity.


\par 
Proper functioning of the internet engineering profession requires certain freedoms and
demands certain societal agreements that are needed to bring about collaborative and
cumulative environments for development of poly-existentials in general and software and
Internet services in particular.


\par Our goals are three fold:

\begin{enumerate}
  \item 
Within the existing corrupt Western legal framework of Copyright and Patents,
facilitate development of 
\href{http://www.freeprotocols.org/}{Free Protocols}
,
\href{http://www.fsf.org/}{Halaal/Libre Software}
and 
\href{http://www.libreservices.org/}{Halaal/Libre Services}
.
\item Work towards Abolishment of the Western so-called Intellectual Property Rights (IPR)
  Regime.  We believe that the Western IPR Regime of Patents and Copyright are an obstacle
  to progress of the engineering profession and are detrimental to society.  In:
\begin{center}
\href{http://www.persoarabic.org/PLPC/120033}{The Nature of Poly-Existentials:}\\
Basis for Abolishment of The Western So-Called Intellectual Property Rights Regime\\
\end{center}
we describe our logic and strategy.
\item Promote development and usage of halaal \begin{fa}«حلال»\end{fa} manner-of-existence
  of software and halaal manner-of-existence of internet application services.  Discourage
  development and usage of Haraam Software and Haraam Internet Application Services.  We
  do this based on an ethical belief system that recognizes responsibility of the
  engineering profession to society, which in turn demands the Halaal manner-of-existence
  of Software. See
\begin{center}
\href{http://www.persoarabic.org/PLPC/120041}{Defining Halaal Manner-Of-Existence Of Software}\\
And Defining Halaal Internet Application Services\\
\end{center}
for more details. Based on this definition, software and Internet
services such as Microsoft Windows, Apple, Google, Facebook, etc, are
haraam because their manner-of-existence is proprietary, restricted
and internally opaque. 
Our use of the sensitive words ``halaal'' and
``haraam'' are described in:
\begin{center}
\href{http://www.persoarabic.org/PLPC/120041}{Introducing Halaal and Haraam into Globish}\\
Based on Moral Philosophy of Abstract Halaal\\
\end{center}
\end{enumerate}

Ultimately, we want to arrive at societal policies that are based on the halaal
manner-of-existence of software.

Towards the above goals, FPF supports a number of projects.  FPF projects model permits
smaller teams to take on specific aspects of our overall activities and move those forward
effectively.  FPF projects include:

\tableofcontents


\section{Halaal Internet Application Services}

Towards the goal of promoting development and use of Halaal Software and Halaal Internet
Application Services within the existing Western legal framework of Copyright and Patents,
we have formed the general frame works of ``Libre Services'', ``Autonomous Libre Services'' and 
``Federations Of Autonomous Services.''

Top level web page for this project is: \href{http://www.freeprotocols.org/projects/libreServices}{http://www.freeprotocols.org/projects/libreServices}.

\subsection{Libre Services (LibreServices.org)}

Libre Services are free software ideology for Internet services.

The main web site for this project is:
\href{http://www.libreservices.org}{libreservices.org}

Key relevant documents for this project are:

\begin{quote}
  \textbf{Libre Services}\\
  \textbf{A non-proprietary model for delivery of Internet services}\\
  \href{http://www.freeprotocols.org/PLPC/100101}{http://www.freeprotocols.org/PLPC/100101}
\end{quote}

\subsection{Autonomous Libre Services}

An Autonomous Libre Services is a Libre Service where the primary user of the service is
the ``owner'' of the service. 
The concept of Autonomous Libre Services focuses on preserving
privacy by providing control over information to users with safe guards of
functional transparency, information portability and non-retention controls.
Formal definitions are provided in:
\begin{quote}
  \textbf{Defining Halaal Manner-Of-Existence Of Software}\\
  \textbf{And Defining Halaal Internet Application Services}\\
  \href{http://www.freeprotocols.org/PLPC/120041}{http://www.freeprotocols.org/PLPC/120041}
\end{quote}

% Web sites relevant to this project are:
% \href{http://www.by-star.net}{by-star.net} 
% \href{http://www.byname.net}{byname.net} 

\subsection{Federated Libre Services}

The concept of Federation of Libre Services is layered above Autonomous Libre Services and
focuses on interactions amongst Autonomous Libre Services and faciltation of information
aggregation amongst Libre Services.

Key relevant documents for this project is:
\begin{quote}
  \textbf{ByStar Federation of Autonomous Libre Services}\\
  \textbf{Overview of Concepts and Models}\\
  \href{http://www.freeprotocols.org/PLPC/180015}{http://www.freeprotocols.org/PLPC/180015}
\end{quote}

\section{Blee (ByStar Libre Emacs Environment)}

Key relevant documents for this project is:

\begin{quote}
  \textbf{Blee -- ByStar Libre Emacs Environment:}\\
  \textbf{A User Environment For The Halaal/Libre ByStar Digital Ecosystem}\\
  \href{http://www.freeprotocols.org/PLPC/180004}{http://www.freeprotocols.org/PLPC/180004}
\end{quote}

\section{Halaal/Libre ByStar Digital Ecosystem}

Key relevant documents for this project is:

\begin{quote}
  \textbf{The Halaal/Libre ByStar Digital Ecosystem}\\
  \textbf{A Moral Alterantive To The Proprietary Digital Ecosystem}\\
  \href{http://www.freeprotocols.org/PLPC/180016}{http://www.freeprotocols.org/PLPC/180016}
\end{quote}

\section{PersoArabic.org}

PersoArabic.org focuses on Halaal/Libre ByStar Digital Ecosystem in the context of
writing, web publishing, communicating and processing in Persian and in Arabic.

The main web site for this project is:
\href{http://www.persoarabic.org}{persoarabic.org}

Top level web page for this project is: \href{http://www.freeprotocols.org/projects/persoArabic}{http://www.freeprotocols.org/projects/persoArabic}.

\section{The Wap Trap}

In 2000, we provided a detailed criticism of a set of specifications called the ``Wireless
Application Protocol'', or WAP. We warned that WAP is booby-trapped with patents and that
WAP is entirely unfit to play the role of a Mobile Messaging industry standard. 

Top level web page for this project is: \href{http://www.freeprotocols.org/projects/wapTrap}{http://www.freeprotocols.org/projects/wapTrap}.

Key relevant documents for this project are:

\begin{quote}
  \textbf{The WAP Trap}\\
  \textbf{An Expose of the Wireless Application Protocol}\\
  \href{http://www.freeprotocols.org/PLPC/100014}{http://www.freeprotocols.org/PLPC/100014}
\end{quote}

%HEVEA\smallskip

\begin{quote}
  \textbf{Le WAP a la trappe}\\
  \textbf{Un expose sur le Wireless Application Protocol}\\
  \href{http://www.freeprotocols.org/PLPC/100015}{http://www.freeprotocols.org/PLPC/100015}
\end{quote}

%HEVEA\smallskip

\begin{quote}
  \textbf{WAP - The Patent Problem}\\
  \href{http://www.freeprotocols.org/PLPC/100025}{http://www.freeprotocols.org/PLPC/100025}
\end{quote}

\section{Libre Texting}

Top level web page for this project is: \href{http://www.freeprotocols.org/projects/libreTexting}{http://www.freeprotocols.org/projects/libreTexting}.

Key relevant documents for this project is:

\begin{quote}
  \textbf{Libre Texting}\\
  \textbf{A Collaborative Initiative And a Reference Implementation}\\
  \href{http://www.freeprotocols.org/PLPC/110015}{http://www.freeprotocols.org/PLPC/110015}
\end{quote}

\section{Patent Battles -- Defense Against Patent Assertions}

Top level web page for this project is: \href{http://www.freeprotocols.org/projects/patentBattles}{http://www.freeprotocols.org/projects/patentBattles}.

% Key relevant documents for this project are: 

% \begin{quote}
%   \textbf{FPF Position Statement July 31, 2001 -- U.S. Patent \# 6,219,694}\\
%   \href{http://www.freeprotocols.org/PLPC/100203}{http://www.freeprotocols.org/PLPC/100203}
% \end{quote}

% %HEVEA\smallskip

% \begin{quote}
%   \textbf{FPF Position Statement September 12, 2002 -- U.S. Patent \# 6,219,694}\\
%   \href{http://www.freeprotocols.org/PLPC/100202}{http://www.freeprotocols.org/PLPC/100202}
% \end{quote}


\section{FPF Patent Free Protocols Procedures}

This project focuses on establishment of the Free Protocols Foundations
processses to ensure the development and maintenance of patent-free protocols.

Key relevant documents for this project are:

\begin{quote}
  \textbf{The Free Protocols Foundation Policies and Procedures}\\
  \href{http://www.freeprotocols.org/PLPC/100201}{http://www.freeprotocols.org/PLPC/100201}
\end{quote}
%HEVEA\smallskip
\begin{quote}
  \textbf{Author's Patent Free Statement for RFC-2188}\\
  \href{http://www.freeprotocols.org/Records/200002071}{http://www.freeprotocols.org/Records/200002071}
\end{quote}
%HEVEA\smallskip
\begin{quote}
  \textbf{Author's Patent Free Statement for RFC-2524}\\
  \href{http://www.freeprotocols.org/Records/200002072}{http://www.freeprotocols.org/Records/200002072}
\end{quote}

\section{Towards Abolishment Of Western IPR Regime}

Key relevant documents for this project is:

\begin{quote}
  \textbf{The Nature of Poly-Existentials:}\\
  \textbf{Basis for Abolishment of  The Western So-Called Intellectual Property Rights Regime}\\
  \href{http://www.freeprotocols.org/PLPC/120033}{http://www.freeprotocols.org/PLPC/120033}
\end{quote}

